\documentclass[10pt, aspectratio=169]{beamer}

% ---------- Theme ----------
\usetheme{Madrid}
\usecolortheme{default}

% ---------- Packages ----------
\usepackage[utf8]{inputenc}
\usepackage[T2A]{fontenc}
\usepackage[russian, english]{babel}
\usepackage{amsmath, amssymb}
\usepackage{booktabs}
\usepackage{array}
\usepackage{listings}
\usepackage{xcolor}

% ---------- Code style ----------
\lstdefinestyle{textstyle}{
    basicstyle=\ttfamily\small,
    backgroundcolor=\color{gray!10},
    frame=single,
    breaklines=true,
    columns=fullflexible
}

\lstset{style=textstyle}

% ---------- Title ----------
\title[VLE в Cassandra]{Моделирование VLE и азеотропной области \\ в системе вода--этанол--н-гексан}
\author{Артём Вотяков}
\institute{Cassandra Monte Carlo}
\date{\today}

\begin{document}

% ============================================================
\begin{frame}
    \titlepage
\end{frame}

% ============================================================
\begin{frame}{Цель проекта}

\textbf{Цель:} исследовать равновесие жидкость--пар в тройной системе
вода--этанол--н-гексан и сравнить результаты Monte Carlo моделирования
с экспериментальными точками.

\vspace{0.4cm}

\textbf{Что ищем:}

\begin{itemize}
    \item состав жидкой фазы:
    \[
        x_i
    \]
    \item состав паровой фазы:
    \[
        y_i
    \]
    \item условие VLE:
    \[
        \mu_i^{liq} = \mu_i^{vap}
    \]
    \item условие азеотропа:
    \[
        x_i \approx y_i
    \]
\end{itemize}

\end{frame}

% ============================================================
\begin{frame}[fragile]{Используемые инструменты}

Для моделирования используется пакет \textbf{Cassandra Monte Carlo}.

\vspace{0.3cm}

\textbf{Метод:}

\begin{lstlisting}
GEMC-NPT = Gibbs Ensemble Monte Carlo
\end{lstlisting}

\vspace{0.2cm}

\textbf{Почему GEMC:}

\begin{itemize}
    \item моделирует две фазы без явной поверхности раздела;
    \item box 1 соответствует жидкой фазе;
    \item box 2 соответствует паровой фазе;
    \item молекулы переносятся между фазами через swap-ходы;
    \item объёмы фаз меняются независимо.
\end{itemize}

\end{frame}

% ============================================================
\begin{frame}[fragile]{Система и компоненты}

\textbf{Моделируемая смесь:}

\vspace{0.3cm}

\begin{table}
\centering
\begin{tabular}{ccc}
\toprule
\textbf{Species} & \textbf{Компонент} & \textbf{Файл модели} \\
\midrule
1 & вода & \texttt{water.mcf} \\
2 & этанол & \texttt{ethanol.mcf} \\
3 & н-гексан & \texttt{hexane.mcf} \\
\bottomrule
\end{tabular}
\end{table}

\vspace{0.3cm}

Для каждой экспериментальной точки из серии жкспериментальных точек задаются жидкие мольные доли:
\end{frame}

% ============================================================
\begin{frame}[fragile]{Начальные условия моделирования}

Для каждой строки CSV создаётся отдельный расчёт.

\vspace{0.25cm}

\begin{table}
\centering
\begin{tabular}{lrrrr}
\toprule
\textbf{Фаза} & \textbf{Вода} & \textbf{Этанол} & \textbf{Гексан} & \textbf{Всего} \\
\midrule
box 1, жидкость & по $x_i$ & по $x_i$ & по $x_i$ & 2400 \\
box 2, пар      & по $x_i$ & по $x_i$ & по $x_i$ & 600 \\
\midrule
Итого & & & & 3000 \\
\bottomrule
\end{tabular}
\end{table}

\vspace{0.25cm}

Газовая фаза получает тот же начальный состав, но меньшее число частиц.

\vspace{0.25cm}

\textbf{Размеры боксов:}

\begin{lstlisting}
box1 = 71.4 Angstrom
box2 = 300.0 Angstrom
\end{lstlisting}

\textbf{Текущие условия:}

\begin{lstlisting}
T = 329.55 K
P = 1.01325 bar
\end{lstlisting}


\end{frame}

% ============================================================
\begin{frame}[fragile]{Потенциалы и силовое поле}

В модели учитываются межмолекулярные взаимодействия:

\begin{lstlisting}
Lennard-Jones 12-6
Coulomb + Ewald summation
Lorentz-Berthelot mixing rule
\end{lstlisting}

\vspace{0.2cm}

\textbf{Lennard-Jones потенциал:}

\[
    U_{LJ}(r)
    =
    4 \varepsilon
    \left[
        \left( \frac{\sigma}{r} \right)^{12}
        -
        \left( \frac{\sigma}{r} \right)^6
    \right]
\]

\textbf{Кулоновский вклад:}

\[
    U_{coul}(r)
    =
    \frac{q_i q_j}{4 \pi \varepsilon_0 r}
\]

\textbf{Правило смешения Lorentz--Berthelot:}

\[
    \sigma_{ij}
    =
    \frac{\sigma_i + \sigma_j}{2},
    \qquad
    \varepsilon_{ij}
    =
    \sqrt{\varepsilon_i \varepsilon_j}
\]

\end{frame}


% ============================================================
\begin{frame}{Тернарная диаграмма}

Траектория состава строится на тернарной диаграмме.

\vspace{0.3cm}

БУДТЕТ ФОТО
\end{frame}

% ============================================================
\begin{frame}[fragile]{Что сравниваем с экспериментом}

Для каждой экспериментальной точки сравниваются:

\begin{lstlisting}
x_exp  vs x_model
y_exp  vs y_model
K_exp  vs K_model
\end{lstlisting}

\vspace{0.3cm}

Основной вопрос:

\[
    \text{сходится ли GEMC-модель к экспериментальному VLE?}
\]

\vspace{0.3cm}

Для азеотропной области дополнительно проверяется:

\[
    |x_i - y_i| \rightarrow 0
\]

\[
    K_i \rightarrow 1
\]

\vspace{0.3cm}

Если модель стабильно отличается от эксперимента, возможные причины:

\begin{itemize}
    \item параметры силового поля;
    \item правило смешения;
    \item недостаточная длина расчёта;
    \item плохая статистика swap/regrowth ходов.
\end{itemize}

\end{frame}


\end{document}